
This study demonstrates that attention entropy over entity spans provides a statistically robust diagnostic signal for distinguishing entity-substitution errors from other failure modes in LLM benchmark evaluation (p < 0.001, Cohen's d = -1.13). Threshold-based classification achieves 86.7\% accuracy on held-out test data, enabling automated failure routing without per-instance human annotation.

The framework automates attention-based failure diagnosis across 73 benchmark failures, exceeding the 10-20 manual example scale typical of prior interpretability work. Zero-entropy entity-errors (60\% of cases) reveal that entity-substitution failures are precision errors --- the model deterministically attends away from the correct entity rather than diffusing attention broadly.

Structural validation of matched correction routing was performed using synthetic experiments with configurable mock success rates. Results demonstrated pipeline viability but did not validate real-world correction effectiveness. Wikipedia API retrieval, LLM generation, and GPT-judge evaluation were not deployed.

Limitations include single-model validation (GPT-2 only), manual gold labeling requirement (N=100), 27\% sample loss due to tokenizer mismatch, single failure-type focus (entity-substitution only), and unvalidated correction effectiveness (synthetic results only).

Future work includes multi-model replication (Llama-2-7B, GPT-3.5, GPT-4), real-world correction validation (Wikipedia API + GPT-judge), automated failure labeling (supervised classifier trained on 100 gold samples), sub-word-aware span alignment (reduce sample loss), and extension to multiple failure types (reasoning errors, knowledge gaps).

The contribution is a validated method for automating interpretability-based failure diagnosis at benchmark scale, demonstrating that attention patterns serve as actionable diagnostic signals for classification without human annotation. The framework integrates benchmark evaluation, interpretability analysis, and targeted correction routing into a systematic pipeline. Real-world deployment of the correction component remains future work.
